\documentclass[11pt]{article}
\usepackage{mathblatt}
% Lösungen nach mathe-nachhilfe gemeinsam.md „Lösungen“ und pruefung.md § 6; Stufenköpfe wortgleich wie im Blatt (Nachzug 12.10.).
\newcommand{\kennung}{P4Z}
\newcommand{\pkt}{\textperiodcentered{}}
\begin{document}
\pfheftstil
\pflaufkopf[P10 \pkt{} Prüfung]{Pythagoras \pkt{} Lösungen}{\kennung}\pfseitevon
\pfheftkopf{Pythagoras}{Lösungen}
\pfstufenzeile{Gleichung aufstellen \pkt{} Hypotenuse gesucht \pkt{} Kathete gesucht \pkt{} in Figuren versteckt}

\pfstufekopf{Gleichung aufstellen}{letzte 5 Jahre: 2026 \pkt{} 2024 \pkt{} 2022}
\begin{pfloesungb}{}
\lzb{1.}{a)}{$k^2 = m^2 + n^2$}{Hypotenuse $k$ (gegenüber dem rechten Winkel)}
\lzb{}{b)}{$r^2 = p^2 + q^2$}{Hypotenuse $r$}
\lzb{}{c)}{$z^2 = s^2 + t^2$}{Hypotenuse $z$}
\end{pfloesungb}
\begin{pfloesungb}{}
\lzb{2.}{}{B}{Kathetenquadrate zusammen $=$ Hypotenusenquadrat}
\end{pfloesungb}
\begin{pfloesungb}{}
\lzb{3.}{}{C}{$w$ ist die Hypotenuse: $w^2 = u^2 + v^2$}
\end{pfloesungb}

\pfstufekopf{Hypotenuse gesucht}{letzte 5 Jahre: 2025}
\begin{pfloesungb}{}
\lzb{4.}{}{a) und c)}{das Fragezeichen liegt dem rechten Winkel gegenüber; in b) ist eine Kathete gesucht}
\end{pfloesungb}
\begin{pfloesungb}{}
\lzb{5.}{a)}{$53$ cm}{$\sqrt{45^2 + 28^2} = \sqrt{2809}$}
\lzb{}{b)}{$39$ cm}{$\sqrt{36^2 + 15^2} = \sqrt{1521}$}
\lzb{}{c)}{$\approx 9{,}8$ cm}{$\sqrt{4^2 + 9^2} = \sqrt{97}$}
\end{pfloesungb}
\begin{pfloesungb}{}
\lzb{6.}{}{$x \approx 200{,}7$ cm}{$x^2 = 200^2 + 17^2 = 40\,289$}
\end{pfloesungb}

\pfstufekopf{Kathete gesucht}{letzte 5 Jahre: 2026 \pkt{} 2024 \pkt{} 2022}
\begin{pfloesungb}{}
\lzb{7.}{}{B}{$m$ ist die Hypotenuse: $k^2 = m^2 - g^2$}
\end{pfloesungb}
\begin{pfloesungb}{}
\lzb{8.}{a)}{$84$ cm}{$\sqrt{85^2 - 13^2} = \sqrt{7056}$}
\lzb{}{b)}{$\approx 9{,}2$ cm}{$\sqrt{11^2 - 6^2} = \sqrt{85}$}
\end{pfloesungb}
\begin{pfloesungb}{}
\lzb{9.}{}{$\overline{FA} \approx 296{,}9$ m}{$\overline{FA}^2 = 396^2 - 262^2 = 88\,172$}
\end{pfloesungb}
\begin{pfloesungb}{}
\lzb{10.}{}{$\overline{AD} \approx 11{,}0$ m}{Dreieck $ADC$: $\overline{AD}^2 = 13{,}2^2 - 7{,}3^2 = 120{,}95$}
\end{pfloesungb}
\begin{pfloesungb}{}
\lzb{11.}{}{$h \approx 34{,}5$ cm}{Dreieck $BFC$: $h^2 = 38^2 - 16^2 = 1188$}
\end{pfloesungb}

\pfstufekopf{in Figuren versteckt}{letzte 5 Jahre: 2026 \pkt{} 2022}
\begin{pfloesungb}{}
\lzb{12.}{a)}{Höhe eingezeichnet}{Dreieck aus halber Grundseite ($6$ cm), Höhe $x$ und Schenkel ($17$ cm); rechter Winkel am Fußpunkt}
\lzb{}{b)}{Dreieck aus zwei Seiten und Diagonale}{$8$ cm, $5$ cm, $x$; rechter Winkel in der Ecke}
\lzb{}{c)}{Dreieck von vorn}{Durchmesser ($6$ cm), Höhe ($9$ cm), Stab $x$; rechter Winkel am Boden}
\end{pfloesungb}
\begin{pfloesungb}{}
\lzb{13.}{}{$h = 63$ cm}{halbe Grundseite $16$ cm: $h^2 = 65^2 - 16^2 = 3969$}
\end{pfloesungb}
\begin{pfloesungb}{}
\lzb{14.}{}{$\approx 12{,}1$ cm}{Durchmesser $6{,}8$ cm; im Becher $\sqrt{7{,}5^2 + 6{,}8^2} = \sqrt{102{,}49} \approx 10{,}1$ cm; dazu $2$ cm}
\end{pfloesungb}
\end{document}
